% Common Preamble for IITM BS Textbook Project
\usepackage[utf8]{inputenc}
\usepackage[T1]{fontenc}
\usepackage{geometry}
\geometry{
    a4paper,
    left=0.9in,
    right=0.9in,
    top=1.0in,
    bottom=1.0in
}

% Typography
\usepackage{palatino} % Clean serif body font
\usepackage[scaled]{helvet} % Helvetica for sans-serif
\usepackage{microtype} % Better word spacing and justifications

% Math Packages
\usepackage{amsmath}
\usepackage{amssymb}
\usepackage{amsthm}
\usepackage{mathtools}
\usepackage{bm} % Bold math symbols
\usepackage{array}
\usepackage{booktabs}
\usepackage{pgfplots}
\pgfplotsset{compat=1.18}
\IfFileExists{venndiagram.sty}{%
  \usepackage{venndiagram}%
}{%
  \makeatletter
  \newenvironment{venndiagram2sets}[1][]{%
    \begin{tikzpicture}[scale=0.55, transform shape]%
      \draw[thick] (-2,-1.5) rectangle (3.5,1.5);
      \node at (3.1,1.1) {$S$};
      \draw[thick] (0,0) circle (1.2cm);
      \draw[thick] (1.5,0) circle (1.2cm);
      \node at (-0.5,0.7) {$A$};
      \node at (2.0,0.7) {$B$};
      \providecommand{\fillA}{\fill[gray!40,opacity=0.6] (0,0) circle (1.2cm);}%
      \providecommand{\fillB}{\fill[gray!40,opacity=0.6] (1.5,0) circle (1.2cm);}%
      \providecommand{\fillACapB}{%
        \begin{scope}%
          \clip (0,0) circle (1.2cm);%
          \fill[gray!60] (1.5,0) circle (1.2cm);%
        \end{scope}%
      }%
      \providecommand{\fillOnlyA}{%
        \begin{scope}%
          \clip (-2,-1.5) rectangle (3.5,1.5);%
          \fill[gray!40] (0,0) circle (1.2cm);%
          \fill[white] (1.5,0) circle (1.2cm);%
          \draw[thick] (0,0) circle (1.2cm);%
          \draw[thick] (1.5,0) circle (1.2cm);%
        \end{scope}%
      }%
      \providecommand{\fillNotA}{%
        \fill[gray!30] (-2,-1.5) rectangle (3.5,1.5);%
        \fill[white] (0,0) circle (1.2cm);%
        \draw[thick] (0,0) circle (1.2cm);%
        \draw[thick] (1.5,0) circle (1.2cm);%
      }%
  }{%
    \end{tikzpicture}%
  }
  \makeatother
}

% Headings styling
\usepackage{titlesec}
\titleformat{\chapter}[display]
  {\normalfont\sffamily\huge\bfseries\color{primarydark}}
  {\chaptername\ \thechapter}{20pt}{\Huge}
\titleformat{\section}
  {\normalfont\sffamily\Large\bfseries\color{primarydark}}
  {\thesection}{1em}{}
\titleformat{\subsection}
  {\normalfont\sffamily\large\bfseries\color{primarydark}}
  {\thesubsection}{1em}{}

% Colors
\usepackage{xcolor}
% Tailored Slate Blue/Navy for Theorems
\definecolor{primarydark}{HTML}{1e293b}
\definecolor{primarylight}{HTML}{eff6ff}
\definecolor{primaryborder}{HTML}{3b82f6}

% Tailored Teal for Definitions
\definecolor{tealdef}{HTML}{0f766e}
\definecolor{tealbg}{HTML}{f0fdfa}
\definecolor{tealborder}{HTML}{14b8a6}

% Tailored Purple for Examples
\definecolor{purpleex}{HTML}{6b21a8}
\definecolor{purplebg}{HTML}{faf5ff}
\definecolor{purpleborder}{HTML}{a855f7}

% Tailored Coral for Key Formulas
\definecolor{coralformula}{HTML}{c2410c}
\definecolor{coralbg}{HTML}{fff7ed}
\definecolor{coralborder}{HTML}{f97316}

% Tailored Amber for Warnings
\definecolor{amberwarning}{HTML}{b45309}
\definecolor{amberbg}{HTML}{fffbeb}
\definecolor{amberborder}{HTML}{f59e0b}

% Tailored Crimson/Ruby for Exam Gotchas & Misconceptions
\definecolor{traprad}{HTML}{991b1b}
\definecolor{trapbg}{HTML}{fef2f2}
\definecolor{trapborder}{HTML}{ef4444}

% Semantic Accent & Forest Colors for TikZ
\definecolor{forest}{HTML}{166534}
\definecolor{accent}{HTML}{2563eb}

% YouTube Dark Slate card
\definecolor{ytcardbg}{HTML}{f8fafc}
\definecolor{ytcardborder}{HTML}{cbd5e1}
\definecolor{yttext}{HTML}{ef4444} % YouTube Red for links

% Indexing, QR Codes, Listings & Lists
\usepackage{imakeidx}
\IfFileExists{qrcode.sty}{
    \usepackage{qrcode}
}{
    \newcommand{\qrcode}[2][]{}
}
\usepackage{listings}
\usepackage[inline]{enumitem}

% Hyperref configuration (loaded BEFORE custom tcolorbox environments using \href)
\usepackage{hyperref}
\hypersetup{
    colorlinks=true,
    linkcolor=blue,
    filecolor=magenta,      
    urlcolor=blue,
    pdftitle={IITM BS Textbook},
    pdfpagemode=FullScreen,
}

% Custom tcolorboxes
\usepackage[most]{tcolorbox}
\usetikzlibrary{shadows, calc, arrows.meta, shapes.geometric, positioning}

% Definition Box
\newtcolorbox{definitionbox}[1]{
    colback=tealbg,
    colframe=tealborder,
    coltitle=white,
    fonttitle=\bfseries\sffamily,
    title={Definition: #1},
    enhanced,
    attach boxed title to top left={yshift=-2mm, xshift=4mm},
    boxed title style={colback=tealdef, colframe=tealborder, arc=1mm},
    arc=2mm,
    boxrule=0.8mm,
    top=4mm,
    bottom=3mm,
    left=3mm,
    right=3mm,
    drop shadow={opacity=0.1, color=black}
}

% Theorem/Rule Box
\newtcolorbox{theorembox}[1]{
    colback=primarylight,
    colframe=primaryborder,
    coltitle=white,
    fonttitle=\bfseries\sffamily,
    title={Theorem/Rule: #1},
    enhanced,
    attach boxed title to top left={yshift=-2mm, xshift=4mm},
    boxed title style={colback=primarydark, colframe=primaryborder, arc=1mm},
    arc=2mm,
    boxrule=0.8mm,
    top=4mm,
    bottom=3mm,
    left=3mm,
    right=3mm,
    drop shadow={opacity=0.1, color=black}
}

% Example Box
\newtcolorbox{examplebox}[1]{
    colback=purplebg,
    colframe=purpleborder,
    coltitle=white,
    fonttitle=\bfseries\sffamily,
    title={Example: #1},
    enhanced,
    attach boxed title to top left={yshift=-2mm, xshift=4mm},
    boxed title style={colback=purpleex, colframe=purpleborder, arc=1mm},
    arc=2mm,
    boxrule=0.8mm,
    top=4mm,
    bottom=3mm,
    left=3mm,
    right=3mm,
    drop shadow={opacity=0.1, color=black}
}

% Intuition Box
\newtcolorbox{intuitionbox}[1]{
    colback=purple!5!white,
    colframe=purple!75!black,
    coltitle=white,
    fonttitle=\bfseries\sffamily,
    title={Core Intuition: #1},
    enhanced,
    attach boxed title to top left={yshift=-2mm, xshift=4mm},
    boxed title style={colback=purple!75!black, colframe=purple!90!black, arc=1mm},
    arc=2mm,
    boxrule=0.8mm,
    top=4mm,
    bottom=3mm,
    left=3mm,
    right=3mm,
    drop shadow={opacity=0.1, color=black}
}

% Key Formula Box
\newtcolorbox{formulabox}[1]{
    colback=coralbg,
    colframe=coralborder,
    coltitle=white,
    fonttitle=\bfseries\sffamily,
    title={Key Formula: #1},
    enhanced,
    attach boxed title to top left={yshift=-2mm, xshift=4mm},
    boxed title style={colback=coralformula, colframe=coralborder, arc=1mm},
    arc=2mm,
    boxrule=0.8mm,
    top=4mm,
    bottom=3mm,
    left=3mm,
    right=3mm,
    drop shadow={opacity=0.1, color=black}
}

% Warning Box (Common Mistakes)
\newtcolorbox{warningbox}[1]{
    colback=amberbg,
    colframe=amberborder,
    coltitle=white,
    fonttitle=\bfseries\sffamily,
    title={Warning: #1},
    enhanced,
    attach boxed title to top left={yshift=-2mm, xshift=4mm},
    boxed title style={colback=amberwarning, colframe=amberborder, arc=1mm},
    arc=2mm,
    boxrule=0.8mm,
    top=4mm,
    bottom=3mm,
    left=3mm,
    right=3mm,
    drop shadow={opacity=0.1, color=black}
}

% Exam Trap & Common Misconception Box (Ruby/Crimson)
\newtcolorbox{trapbox}[1]{
    colback=trapbg,
    colframe=trapborder,
    coltitle=white,
    fonttitle=\bfseries\sffamily,
    title={Exam Trap \& Common Misconception: #1},
    enhanced,
    attach boxed title to top left={yshift=-2mm, xshift=4mm},
    boxed title style={colback=traprad, colframe=trapborder, arc=1mm},
    arc=2mm,
    boxrule=0.8mm,
    top=4mm,
    bottom=3mm,
    left=3mm,
    right=3mm,
    drop shadow={opacity=0.1, color=black}
}

% YouTube Video Resource Card
\newtcolorbox{videocard}[2]{
    colback=ytcardbg,
    colframe=ytcardborder,
    enhanced,
    boxrule=0.5mm,
    arc=1.5mm,
    left=3mm,
    right=3mm,
    top=3mm,
    bottom=3mm,
    fontupper=\small\sffamily,
    title={\href{#2}{\textcolor{yttext}{\textbf{Video Resource:}} \textcolor{primarydark}{#1}} \hfill \raisebox{-1.5mm}{\qrcode[height=0.6cm]{#2}}},
    coltitle=primarydark,
    colbacktitle=ytcardbg,
    attach boxed title to top left={yshift=-2mm, xshift=2mm},
    boxed title style={colback=ytcardbg, colframe=ytcardbg},
    underlay={
        \path[draw=ytcardborder, line width=0.5mm] (frame.north west) -- (frame.south west);
    },
    title style={bottom=1mm}
}

% Python Code Snippets Styling
\definecolor{codegreen}{HTML}{10B981}
\definecolor{codegray}{HTML}{6B7280}
\definecolor{codepurple}{HTML}{8B5CF6}
\definecolor{codebg}{HTML}{F8FAFC}
\lstdefinestyle{pystyle}{
    language=Python,
    backgroundcolor=\color{codebg},   
    commentstyle=\color{codegray}\itshape,
    keywordstyle=\color{codepurple}\bfseries,
    numberstyle=\tiny\color{codegray},
    stringstyle=\color{codegreen},
    basicstyle=\ttfamily\footnotesize,
    breakatwhitespace=false,         
    breaklines=true,                 
    captionpos=b,                    
    keepspaces=true,                 
    numbers=left,                    
    numbersep=5pt,                  
    showspaces=false,                
    showstringspaces=false,
    showtabs=false,                  
    tabsize=4,
    frame=single,
    rulecolor=\color{ytcardborder},
    xleftmargin=10pt,
    xrightmargin=10pt
}
\lstset{style=pystyle}

% Custom Math Macros (Common across all courses)
\newcommand{\N}{\mathbb{N}}
\newcommand{\Z}{\mathbb{Z}}
\newcommand{\Q}{\mathbb{Q}}
\newcommand{\R}{\mathbb{R}}
\newcommand{\C}{\mathbb{C}}
\newcommand{\set}[1]{\left\{#1\right\}}
\newcommand{\abs}[1]{\left|#1\right|}
\newcommand{\func}[3]{#1 \colon #2 \to #3}
\newcommand{\limit}[2]{\lim_{#1 \to #2}}
\newcommand{\deriv}[2]{\frac{d#1}{d#2}}
\newcommand{\pderiv}[2]{\frac{\partial#1}{\partial#2}}
\newcommand{\dx}{\,dx}
